\documentclass[12pt,a4paper]{article}
\usepackage[utf8]{inputenc}
\usepackage[T1]{fontenc}
\usepackage{enumitem}
\usepackage[margin=2.5cm]{geometry}
\usepackage{parskip}
\usepackage{amsmath}
\usepackage{xcolor}

\title{thesis — Final Exam --- Answer Key}
\date{}

\begin{document}

\maketitle

{\large \textbf{\textcolor{red}{ANSWER KEY --- Not for Distribution}}}

\vspace{0.5em}



\noindent
\textbf{Total points:} 7
\hfill
\textbf{Number of questions:} 7

\vspace{1em}
\hrule
\vspace{1.5em}


\noindent
\textbf{Question 1 (1.00 pts).}~
Which of the following is a characteristic of the first generation of computers?


\begin{enumerate}[label=\textbf{\Alph*.}]

  \item \textbf{\textcolor{teal}{[CORRECT]}}  They were built with vacuum tubes

  \item They were built with superconductors

  \item They were built with artificial intelligence

  \item They were built with advanced computing

\end{enumerate}



\textit{\small Explanation: The first generation of computers were built with vacuum tubes, which were about the size of a light bulb.}



\vspace{0.8em}
\noindent\rule{\linewidth}{0.2pt}
\vspace{0.3em}


\noindent
\textbf{Question 2 (1.00 pts).}~
What does the term 'Epoch' refer to in the context of Neural Networks?


\begin{enumerate}[label=\textbf{\Alph*.}]

  \item A single set of input/output data

  \item \textbf{\textcolor{teal}{[CORRECT]}}  A complete set of input/output data

  \item A type of neural network topology

  \item A type of neural network architecture

\end{enumerate}



\textit{\small Explanation: The term 'Epoch' is defined in the context as a complete set of input/output data.}



\vspace{0.8em}
\noindent\rule{\linewidth}{0.2pt}
\vspace{0.3em}


\noindent
\textbf{Question 3 (1.00 pts).}~
A neural network with two hidden layers is sufficient to solve any problem.


\textbf{Answer:} \textcolor{teal}{True}



\textit{\small Explanation: Two hidden layers are sufficient to solve any problem, as stated in the context.}



\vspace{0.8em}
\noindent\rule{\linewidth}{0.2pt}
\vspace{0.3em}


\noindent
\textbf{Question 4 (1.00 pts).}~
How do the strength and type of synapses affect the firing of a neuron, and what are the implications for artificial neural networks?


\textbf{Model Answer:} \textcolor{teal}{The strength and type of synapses determine the signal strength transmitted to the neuron. Excitatory synapses stimulate the neuron, while inhibitory synapses restrict it. In artificial neural networks, this is modeled using weighted inputs and thresholds.}



\textit{\small Explanation: Synapses vary in strength and type; Excitatory and inhibitory synapses have different effects on neuron firing; Artificial neural networks model this using weighted inputs and thresholds}



\vspace{0.8em}
\noindent\rule{\linewidth}{0.2pt}
\vspace{0.3em}


\noindent
\textbf{Question 5 (1.00 pts).}~
How does the use of neural networks in the context of wireless network reliability analysis differ from traditional reliability analysis methods?


\textbf{Model Answer:} \textcolor{teal}{Neural networks can be trained with results from wireless simulations to improve reliability or maintainability, whereas traditional methods may rely on statistical models or discrete-event simulations alone.}



\textit{\small Explanation: Neural networks can learn from simulation data; Traditional methods may not account for complex interactions}



\vspace{0.8em}
\noindent\rule{\linewidth}{0.2pt}
\vspace{0.3em}


\noindent
\textbf{Question 6 (1.00 pts).}~
What makes neural networks unique and how do they work?


\textbf{Key Points:} \textcolor{teal}{I. Introduction to neural networks, II. Key characteristics of neural networks (nonlinearity, learning from data), III. Training a neural network (introducing data, computing output, modifying weights), IV. Using a neural network (introducing new data, computing output)}



\textit{\small Explanation: Explain the key characteristics of neural networks, including their nonlinearity and ability to learn from data, and describe the process of training and using a neural network.}



\vspace{0.8em}
\noindent\rule{\linewidth}{0.2pt}
\vspace{0.3em}


\noindent
\textbf{Question 7 (1.00 pts).}~
Explain how neural networks can be used to extract knowledge and provide rules that are human comprehensible, and discuss the significance of knowledge extraction in complex processes.


\textbf{Key Points:} \textcolor{teal}{I. Introduction to neural networks and knowledge extraction, II. Significance of knowledge extraction in complex processes, III. Examples of neural networks in knowledge extraction, IV. Conclusion}



\textit{\small Explanation: Consider the relationship between neural networks, knowledge extraction, and complex processes, and provide examples from the context to support your explanation.}





\end{document}
